The JSR was introduced by Rota and Strang~\cite{rota1960note}
and has many other applications such as wavelets, the capacity of some particular codes, zero-order stability of ordinary differential equations, congestion control in computer networks, curve design and networked and delayed control systems; see \cite{jungers2009joint} for a survey on the JSR and its applications.
Many algorithms exist for estimating the JSR but not much is known on how to generate an infinite
sequence of matrices with an asymptotic growth rate close to the JSR.
However generating such sequence can be of particular interest, depending on the application,
such as exhibiting unstable trajectories for s\witch{} linear systems to prevent them from occuring~\cite{gomes2018minimally}.%
The currently known algorithms generate a sequence of matrices with high spectral radius
using methods detailed in \secref{lb}
and repeat this sequence infinitely \cite{gripenberg1996computing,guglielmi2008algorithm,guglielmi2013exact,jungers2014lifted}.

Approximating the JSR usually consists in certifying upper bounds $\gamub$ to the JSR by exhibiting a Lyapunov function or an invariant set for the matrices $A_i/\gamub$.
The search for such Lyapunov functions can naturally be written as a convex optimization program; see \progref{primalinf}.
Certifying lower bounds $\gamlb$ is currently either achieved using the guarantees we have on the accuracy of the upper bound on the JSR
or by exhibiting trajectories of asymptotic growth rate $\gamlb$.
In this paper, we introduce a new way to certify lower bounds by exhibiting nonnegative measures satisfying some invariance condition parametrized by the matrices $A_i/\gamlb$; see \eqref{eq:dualinf1}.
This invariance condition is linear on the measure hence the search of measures on the convex cone of nonnegative measures is a \emph{convex} program; see \progref{dualinf}.
It turns out that this program is the dual of \progref{primalinf}.

We revisit the sum-of-squares program proposed by Parrilo and Jadbabaie~\cite{parrilo2008approximation} and show that its dual formulation is the moment relaxation of the search of the measures satisfying the invariance condition.

Thanks to this duality, solving this pair of programs with a given candidate value $\gamma$
for the JSR either returns Lyapunov functions certifying that
$\jsr \leq \gamma$ or returns moments that form a solution of the moment relaxation; see \secref{moments}.
These moments are not necessarily the moments of measures satisfying the invariance conditions.
However, we give a rounding procedure to extract a (infinite) switching sequence from these moments
and provide a guarantee on the asymptotic growth rate of this sequence.
As a by-product of the rounding procedures, the spectral radius of a finite part of
this infinite sequence can be used to give lower bounds on the JSR.
In addition, we give a way to sometimes detect when the moments belongs to measures that satisfy the invariance conditions.
This happens when the measures are the convex combination of the occupation measures of several periodic trajectories.
Since the trajectories are periodic, the measures are atomic and we can recover them from moments of sufficiently high degree.
We show on numerical examples that these techniques work well in practice.
